\section{Introduction}
\label{sec:intro}

% What is compositionality?
Scholars today herald \emph{compositionality} as a fundamental presupposition characterizing both human perception and linguistic processing~\cite{cresswell1973logics}.
Through compositional reasoning, humans can comprehend new scenes and describe those scenes by composing known atoms~\cite{janssen1997compositionality,hupkes2020compositionality,bottou2014machine,chomsky1965some}.
For instance, compositionality allows people to differentiate between a photo of ``a girl in white facing a man in black'' and ``a girl in black facing a man in white''.
For a while now, vision-language research has sought to develop models that can similarly comprehend scenes and express them through compositional language~\cite{krishna2017visual,ji2020action,lu2016visual,GrundeMcLaughlin2021AGQA}.
% Inspired from Cognitive Science~\cite{biederman1987recognition} today's vision research identifies compositional atoms to consist of objects (e.g.~``vase'', ``television''), their attributes (e.g.~``yellow'', ``black''), and their pairwise relationships (e.g.~``on top of'',)~\cite{krishna2017visual,ji2020action,lu2016visual,schuster2015generating,GrundeMcLaughlin2021AGQA}.


%  What benchmarks have been proposed thus far? How do they evaluate?
Given its importance, a surge of new benchmarks have been proposed to evaluate whether vision-language models exhibit compositionality.
Recently,  Winoground~\cite{thrush2022winoground}, VL-CheckList~\cite{zhao2022vl}, ARO~\cite{yuksekgonul2023when}, CREPE~\cite{ma2022crepe}, and Cola~\cite{ray2023cola} have entered the machine learning zeitgeist.
Evaluation is mostly done through an image-to-text retrieval task formulation: by measuring how often models pick the description, ``a girl in white facing a man in black'' when presented with an image of it, and avoid choosing the incorrect \emph{hard negative} description, ``a girl in black facing a man in white''.


% Problem: blind models outperform others.
In this work, we uncover a crucial vulnerability in not just one but all these image-to-text compositionality benchmarks:
We find that a \emph{blind} model that never looks at the image, can identify the correct caption and avoid choosing the supposed ``hard negatives''.
This blind model outperforms a wide array of pretrained vision-language models across the suite of benchmarks~\cite{RadfordKHRGASAM21,ilharco_gabriel_2021_5143773,gadre2023datacomp}.
We explain this undesired hackability in existing benchmarks by showcasing that there exists a significant distributional gap between the positive and hard negative captions.
For instance, in the ARO benchmark~\cite{yuksekgonul2023when}, human-generated positive captions differ drastically from the hard negative texts generated by randomly shuffling words in the positive captions.
As new research has begun to propose methods that claim to improve compositionality on these benchmarks~\cite{yuksekgonul2023when,ray2023cola}, we find it critical to highlight our findings and propose a solution.


% solution: We present SUGAR-CREPE. How is this constructed. Why does it not have the above flaws ?
We propose a solution to existing hackable benchmarks by introducing \method, a new benchmark to faithfully evaluate compositionality. In curating \method, we identify two main \emph{biases}~\footnote{\ We use biases and artifacts interchangeably in the paper.} that result in the distributional gap between positive and hard negatives; and employ mechanisms to fix the shifts.
In particular, we find the current procedure in generating hard negatives introduces descriptions that are (1) not plausible and (2) non-fluent.
For example, while the caption ``olives and grapes on a plate'' is a sensical fluent caption, benchmarks often have non-plausible hard negatives like ``olives and grapes inside a plate'' or simply incomprehensible ones like ``right has word another word. There is a words'' (see Table~\ref{tab:bad_hn_examples} for more examples).
We mitigate such biases by first leveraging a modern large language model, ChatGPT~\cite{chatgpt}, to generate plausible and natural hard negative texts instead of relying on simple rule-based templates employed by existing benchmarks~\cite{ma2022crepe,yuksekgonul2023when}.
Then, we subsample the dataset through an adversarial refinement process to ensure the identified biases are maximally removed by drawing on recent dataset de-biasing work~\cite{zellers2018swag,sakaguchi2021winogrande,le2020adversarial}.
Taken together, this workflow is where \method derived its name: \textbf{S}ynthetic yet \textbf{U}nbiased \textbf{G}eneration with \textbf{A}dversarially \textbf{R}efined \textbf{C}ompositional \textbf{REP}resentation \textbf{E}valuation.
We qualitatively and quantitatively verify through both human and automatic evaluations that \method\ effectively fixes these biases.


% Given this new benchmark, what are your new findings when you measure SOTA representations ? Do your findings match the previous ? Or do we have new ones ? What is the implication of SUGAR-CREPE.
With \method, we \emph{re}-evaluate recent methods proposed to improve compositionality. Specifically, we focus on one prominent approach that aims to improve compositionality through data augmentation. This method trains models by generating compositional hard negatives and injecting them within a training batch~\cite{doveh2023teaching,yuksekgonul2023when}.
Unfortunately, we observe that the effectiveness of this simple data augmentation approach is hugely \emph{overestimated} when evaluated on existing benchmarks, leading to limited improvements on \method.
Finally, we evaluate a wide variety of $17$ pretrained CLIP models~\cite{RadfordKHRGASAM21,ilharco_gabriel_2021_5143773,gadre2023datacomp}, and find that current models still lack compositionality.
Our results suggest that to improve compositionality, future work may need more innovative techniques.


% For instance, to test a VLM's capability of understanding atomic concepts, \cite{zhao2022vl,ma2022crepe} generate hard negatives by replacing some words in the original text with others.
% Similarly, hard negatives generated by swapping or shuffling some words in the original text can be used to test a model's ability in semantic understanding~\cite{thrush2022winoground,yuksekgonul2023when,ma2022crepe}.
% As pointed out in these works, current VLMs unfortunately lack compositionality and are specifically insensitive to small atomic or ordering changes~\cite{thrush2022winoground,zhao2022vl,ma2022crepe,yuksekgonul2023when}.
